% Generated draft: requires mathematical and source review.
\documentclass[11pt]{article}
\usepackage[margin=25mm]{geometry}
\usepackage{fontspec}
\setmainfont{Latin Modern Roman}
\usepackage{amsmath,amssymb,mathtools,unicode-math}
\setmathfont{Latin Modern Math}
\usepackage{enumitem,needspace,xurl,hyperref}
\setlength{\parindent}{0pt}
\setlength{\parskip}{5pt}
\setlength{\emergencystretch}{3em}
\newcommand{\N}{\mathbb{N}}
\newcommand{\Z}{\mathbb{Z}}
\newcommand{\Q}{\mathbb{Q}}
\newcommand{\R}{\mathbb{R}}
\newcommand{\C}{\mathbb{C}}
\newcommand{\F}{\mathbb{F}}
\newcommand{\A}{\mathbb{A}}
\newcommand{\PP}{\mathbb{P}}
\newcommand{\E}{\mathbb{E}}
\newcommand{\eps}{\varepsilon}
\newcommand{\dd}{\,\mathrm{d}}
\begin{document}
\section*{UnsolvedMath ID: 30003159}
\textbf{Source problem code:} catalogue entry
\subsection{Definitions and mathematical statement}
A smooth complex projective variety \(V\) of dimension \(n\) is called rational if it is birational over \(\C\) to \(\PP^n_{\C}\). Here birational means that dense Zariski open subsets of the two varieties are isomorphic. It is called irrational if it is not rational.

In the algebraic interpretation of the supplied question, let \(B\) be a nonempty irreducible complex algebraic variety and let \(f\colon X\to B\) be a smooth projective morphism whose fibers \(X_b=f^{-1}(b)\), for \(b\in B(\C)\), are nonempty integral varieties of a common dimension \(n\geq 0\). Is the following assertion true for every such family?
\[
\forall b,c\in B(\C),\qquad
X_b\text{ is rational}\ \Longleftrightarrow\ X_c\text{ is rational}.
\]
Equivalently, does there exist a family satisfying these hypotheses and points \(b,c\in B(\C)\) such that \(X_b\) is rational and \(X_c\) is irrational? Existence of such a family would give a negative answer to deformation invariance. The supplied background also singles out the restriction \(n=3\).

The source does not explicitly specify whether the base is algebraic or analytic, or state an integrality or connectedness convention for fibers. The displayed formulation uses algebraic bases and the usual integral meaning of variety; these conventions require confirmation against the original problem. No smoothness assumption on \(B\) is supplied.
\subsection{Short English statement}
Must rationality be constant among the fibers of every smooth projective family over an irreducible complex base? Equivalently, can such a family contain both a rational fiber and an irrational fiber? The supplied background highlights threefolds. The algebraic-base and integral-fiber conventions need source confirmation.
\subsection{Sources}
\begin{itemize}
\item Local UnsolvedMath registry, permanent ID \(30003159\), title ``Deformation Invariance of Rationality,'' source problem code \texttt{catalogue entry}, in the Oberwolfach Reports Open Problems collection. The registry records proposed year 2016 and supplies no named proposer.
\item Supplied source citation: \emph{Classical Algebraic Geometry (2017)}, Oberwolfach Reports, DOI \texttt{10.4171/owr/2016/29}, pp.\ 1635--1682. Supplied source URL and literature reference ``Oberwolfach Report 29'': \url{https://doi.org/10.4171/owr/2016/29}. The background describes these reports as documenting mathematical workshops and their open-problem discussions.
\item Literature assessment checked on \texttt{2026-08-21}: supplied dated curation, not independently verified here. It labels the problem open and states that rationality is deformation invariant in several important classes, while the general question, including the highlighted dimension-three case, remains open. The supplied curation expressly requests independent verification of time-sensitive status and cited claims.
\item The registry labels the statement \texttt{corrected\_verified} and reports a manual semantic, typographic, and LaTeX review. These are supplied editorial assertions; they do not constitute independent mathematical or bibliographic verification of this draft.
\end{itemize}
\end{document}
